\documentclass[12pt]{article}
\usepackage[pdftex,pagebackref,letterpaper=true,colorlinks=true,pdfpagemode=none,urlcolor=blue,linkcolor=blue,citecolor=blue,pdfstartview=FitH]{hyperref}

\usepackage{amsmath,amsfonts}
\usepackage{graphicx}
\usepackage{color}


\setlength{\oddsidemargin}{0pt}
\setlength{\evensidemargin}{0pt}
\setlength{\textwidth}{6.0in}
\setlength{\topmargin}{0in}
\setlength{\textheight}{8.5in}

\setlength{\parindent}{0in}
\setlength{\parskip}{5px}

\input{macrosblog}

\begin{document}

% A digestion of the proof of Sendov's conjecture

\begin{conjecture}[Sendov's conjecture]  Let $n \geq 2$, and let $p : \C \to \C$ be a degree $n$ polynomial with all zeroes in the unit disk.  Then for every zero $a$ of $p$, there exists a critical point $\zeta$ of $p$ with $|\zeta-a| \leq 1$.
\end{conjecture}

\begin{conjecture}[Phelps--Rodriguez conjecture]  Let $n \geq 2$, and let $p : \C \to \C$ be a degree $n$ polynomial with all zeroes in the unit disk.  Then for every zero $a$ of $p$, there exists a critical point $\zeta$ of $p$ with $|\zeta-a| < 1$, unless $a$ is on the unit circle and $p$ is a scalar multiple of $z^n - a^n$.
\end{conjecture}

By applying a rotation around the origin, we can normalize $a$ to be a real number with $0 \leq a \leq 1$.  

From the \href{???}{work of Rubinstein}, both conjectures were already established in the $a=1$ case, so one can restrict to the $0 \leq a < 1$ case.  Both of these conjectures then follow from

\begin{conjecture}[Sendov's conjecture in interior]\label{interior}  Let $n \geq 2$.  Let $p : \C \to \C$ be a degree $n$ polynomial with all zeroes in the unit disk.  Then if $0 \leq a < 1$ is a zero of $p$, there exists a critical point $\zeta$ of $p$ with $|\zeta-a| < 1$.
\end{conjecture}

All three of these conjectures were established for $n \leq 8$ (in a sequence of papers culminating in \href{???}{this paper of Brown and Xiang}) and for sufficiently large $n$ (in \href{???}{a paper of myself}, which in turn built upon several partial results in this setting).  This left the case of intermediate $n$ to be settled.  Recently, Lech Mazur \href{???}{was able to use an AI tool} to resolve Sendov's conjecture for all $n \geq 2$, with the proof \href{???}{verified in Lean}.  However, the AI-generated proof was not human-digested to be in the form of a publication-ready preprint; and it has taken me several days (with \href{???}{heavy AI assistance}) to perform such a digestion, to place the proof in proper context with previous literature and to simplify and streamline the argument to highlight the main ideas.

One consequence of this digestion is that the argument in fact demonstrates Conjecture \ref{interior}, and thus resolves both the Sendov conjecture and the Phelps--Rodriguez conjecture in full generality.  

For expositional reasons, we will restrict attention here to the case where $n$ is sufficiently large, as one can then handle certain lower order terms cleanly using asymptotic notation.  However, the arguments here extend (after some computer assistance) to cover the range $n \geq 5$, and in fact with additional modifications can also treat the very low degree cases $n=2,3,4$, although these cases were already known in the literature.  I used an AI to provide a full treatment of the $n \geq 5$ case along the lines of this blog post; this can be found \href{???}{here}.

We now prove Conjecture \ref{interior} in the high degree case.  Suppose for contradiction that the claim failed for some sufficiently large $n$, thus one can find a degree $n$ polynomial $p$ with zeroes
$$ a, z_1, \dots,z_{n-1}$$
for some $0 < a  < 1$ and $z_j$, $j=1,\dots,n-1$ in the closed unit disk, whose critical points all lie a distance at least $1$ from $a$.  We use $O()$ notation here in the non-asymptotic sense, thus $X = O(Y)$ means that $|X| \leq C Y$ for some absolute constant $C$ (independent of $n$).  We will also use the notation $O_{\leq}(Y)$ to denote a quantity that is bounded in magnitude by $Y$.

To capture the fact that the critical points lie are at a distance at least $1$ from $a$, we write these critical points as
$$ a - \frac{1}{q_1}, \dots, a - \frac{1}{q_{n-1}}$$
for some (non-zero) $q_1,\dots,q_{n-1}$ in the closed unit disk.

\begin{example}\label{ex}  If $p(z) = z^n - 1$ and $a=1$, then $z_1,\dots,z_{n-1}$ are the non-trivial $n^{th}$ roots of unity, while the $q_1,\dots,q_{n-1}$ are all equal to $1$.  Strictly speaking this is not actually a counterexample to Theorem \ref{interior}, because $a$ is not strictly less than one; nevertheless this is an important motivating near-counterexample for the arguments below.
\end{example}

\begin{example}\label{ex2}  A generalization of the previous example was studied in Section 4 of \href{???}{my paper}.  Here one took
    $$ p(z) = (z + \frac{c_2}{n})^{n-m} P(z) - (a + \frac{c_2}{n})^{n-m} P(a)$$
where $n$ was an asymptotic parameter going to infinity,
$$ P(z) = (z-\lambda_1) \dots (z - \lambda_m)$$
was a low-degree polynomial for some $m=O(1)$,
$$ a = 1 - \frac{c_1}{n},$$
and $c_1,c_2 > 0$ were constants.  This polynomial has a zero at $a$, $n-m-1$ critical points at $-\frac{c_2}{n}$, and $m$ additional critical points near $\lambda_1,\dots,\lambda_m$.  If all the critical points were at distance at least one from $a$, one would have
$$ c_2 \geq c_1$$
and
$$ |1-\lambda_j| \geq 1 - o(1)$$
while if all the zeroes were in the unit disk, the calculations in my paper showed that
\begin{equation}\label{c21} c_2 - c_1 - c_2 \cos \theta + \sum_{j=1}^m \log |\frac{1-\lambda_j}{e^{i\theta}-\lambda_j}| \leq 1 + o(1)
\end{equation}
Here $o(1)$ denotes a quantity that goes to zero as $n \to \infty$.  If one ignores the $o(1)$ errors, one can show that these conditions are only simultaneously feasible if $c_1=c_2$ and all the $\lambda_j$ vanish, but the argument was somewhat subtle (I had to proceed by inspecting the second Fourier coefficient of \eqref{c21}).  This illustrates the fact that the regime $a = 1 - O(1/n)$ is particularly delicate.
\end{example}

We now have two sets of points in the closed unit disk: $z_1,\dots,z_{n-1}$ and $q_1,\dots,q_{n-1}$.  They ``communicate'' with each other through the polynomial $p$ and its first derivative $p'$, both of which can be expressed in terms of either set of points (as well as $n$ and $a$).  Indeed, if we normalize $p$ to be monic, then we can factor $p$ in terms of the zeroes as
\begin{equation}\label{p-eq-zero}
    p(z) = (z-a) \prod_{j=1}^{n-1} (z-z_j)
\end{equation}
and thus upon differentiating
\begin{equation}\label{p-deriv-zero}
    p'(z) = (\prod_{j=1}^{n-1} (z-z_j)) (1 + (z-a) \sum_{j=1}^{n-1} \frac{1}{z-z_j}).
\end{equation}
Here and in the sequel we adopt the convention of removing singularities when dealing with expressions that involve multiplication by both $\frac{1}{z-z_j}$ and $z-z_j$, by cancelling such terms first in the event that $z=z_j$.

In a similar vein, $p'$ can be factored 
\begin{equation}\label{p-factor}
    p'(z) = n \prod_{j=1}^{n-1} (z - a + \frac{1}{q_j})
\end{equation}
and thus on integrating (and using $p(a)=0$)
\begin{equation}\label{p-integral}
    p(z) = (z-a) \int_0^1 n \prod_{j=1}^{n-1} (t(z - a) + \frac{1}{q_j})\ dt.
\end{equation}

It is convenient to rule out the easy case $a=0$ right away.  In this case we see from \eqref{p-deriv-zero}, \eqref{p-factor} that
$$ p'(0) = \prod_{j=1}^{n-1} (-z_j) = n \prod_{j=1}^{n-1} \frac{1}{q_j}$$
which is absurd since the first product has magnitude at most one, and the second product has magnitude at least one.  Thus we can assume henceforth that $a>0$.

By inspecting $p$ or $p'$ at various natural locations, we can thus obtain a number of identities relating the $z_j$ to the $q_j$.  We record the ones that we actually need here:

\begin{lemma}\label{identities}  Let $F$ denote the function
\begin{equation}\label{F-def} F(t) := \prod_{j=1}^{n-1} (1 - atq_j).
\end{equation}
\begin{itemize}
    \item [(i)] (Centroid identity)  We have
    $$ \frac{1}{n} (a + \sum_{j=1}^{n-1} z_j) = \frac{1}{n-1} \sum_{j=1}^{n-1} (a - \frac{1}{q_j}).$$
    That is to say, the centroid of the zeroes equals the centroid of the critical values.
    \item [(ii)]  (Polar identity)  We have
    $$ \prod_{j=1}^{n-1} \frac{1-az_j}{a-z_j} = \int_0^1 \prod_{j=1}^{n-1} (t (1-a^2) q_j + a)\ dt.$$
    \item [(iii)]  (First origin identity) We have
\begin{equation}\label{first-origin} 
    (-1)^{n-1} \prod_{j=1}^{n-1} z_j = \frac{n}{\prod_{j=1}^{n-1} q_j} \int_0^1 F(t)\ dt.
\end{equation}
 \item [(iv)]  (Second origin identity) We have
    $$ (-1)^{n-1} \prod_{j=1}^{n-1} z_j \left( 1 + a \sum_{j=1}^{n-1} \frac{1}{z_j} \right) = \frac{n}{\prod_{j=1}^{n-1} q_j} F(1).$$ 
    (Again, we are using the convention of removing singularities to deal with the case where some of the $z_j$ vanish.)
\end{itemize}
\end{lemma}

\begin{proof}  For (i), we inspect the behavior of $p'(z)$ as $z \to \infty$.  From \eqref{p-eq-zero} we have
    $$ p(z) = z^n - ( a + \sum_{j=1}^{n-1} z_j )  z^{n-1} + O(z^{n-2})$$
    and thus on differentiating term by term
    $$ p'(z) = n z^{n-1} - (n-1) ( a + \sum_{j=1}^{n-1} z_j ) z^{n-2} + O(z^{n-3}).$$
    Meanwhile, from \eqref{p-factor} we have
    $$ p'(z) = n z^{n-1} - n (\sum_{j=1}^{n-1} (a - \frac{1}{q_j})) z^{n-2} + O(z^{n-3})$$
comparing coefficients, we obtain the claim.

For (ii), we consider the expression $p(1/a) / p'(a)$.  On the one hand, from \eqref{p-eq-zero}, \eqref{p-deriv-zero} one has
$$ \frac{p(1/a)}{p'(a)} = \frac{(1/a - a) \prod_{j=1}^{n-1} (1/a - z_j)}{\prod_{j=1}^{n-1} (a-z_j)}.$$
(Note from hypothesis that $a$ cannot be a critical point, so the denominator is non-zero.)
On the other hand, from \eqref{p-factor}, \eqref{p-integral} one has
$$ \frac{p(1/a)}{p'(a)} = \frac{(1/a - a) \int_0^1 n \prod_{j=1}^{n-1} (t(\frac{1}{a} - a) + \frac{1}{q_j})\ dt}{n \prod_{j=1}^{n-1} \frac{1}{q_j}}.$$
Equating the two identities, we obtain (ii) after some algebra.

For (iii), we evaluate $p(0)$.  From \eqref{p-eq-zero} we have
$$     p(0) = -a (-1)^{n-1} \prod_{j=1}^{n-1} z_j$$
while from \eqref{p-integral} we have
$$ p(0) = -a \int_0^1 n \prod_{j=1}^{n-1} (-at + \frac{1}{q_j})\ dt.$$
Equating the two identities, we obtain (iii) after some algebra using \eqref{F-def}.

For (iv), we similarly evaluate $p'(0)$.  From \eqref{p-deriv-zero} we have
$$ p'(0) = (-1)^{n-1} (\prod_{j=1}^{n-1} z_j) (1 + a \sum_{j=1}^{n-1} \frac{1}{z_j})$$
while from \eqref{p-factor} one has
$$ p'(0) = n \prod_{j=1}^{n-1} (- a + \frac{1}{q_j}).$$
Equating the two identities, we obtain (iv) after some algebra using \eqref{F-def}.
\end{proof}

Remarkably, the polynomial $p$ will play no further role in the argument: the identities in (i)-(iv), together with the hypotheses that $0 \leq a < 1$ and $z_j, q_j$ lie in the closed unit disk, will be sufficient by themselves to obtain a contradiction.

\begin{example}  Continuing the example in Example \ref{ex}, in (i) both sides vanish.  In (ii), both sides are equal to one.  For (iii) and (iv), we have $F(t) = (1-t)^{n-1}$, with both sides equal of (iii) equal to one, and both sides of (iv) equal to zero.
\end{example}

\begin{remark} The comparison of the polynomial at a location $a$ and at the inversion $1/a$ of that location across the closed unit disk is a familiar trick in the literature; see, e.g., ???
\end{remark}

\begin{remark} The first origin identity \eqref{first-origin} is already strong enough to handle asymptotically all examples of the form inn Example \ref{1a}, except in the endpoint case where $c_1, c_2$ vanish and the $\lambda_j$ are all $1$.  Indeed, as the $z_j, q_j$ are in the closed unit disk, \eqref{first-origin} implies that
    $$ n |\int_0^1 F(t)\ dt| \leq 1.$$
On the other hand, routine calculations (omitted here) show that
$$ n \int_0^1 F(t)\ dt = 1 + \frac{c_2 + \sum_{j=1}^{n-1} \frac{\lambda_j}{1-\lambda_j}}{n}
+ O(\frac{1}{n^2})$$
leading asymptotically to the constraint
$$ c_2 + \mathrm{Re} \sum_{j=1}^{n-1} \frac{\lambda_j}{1-\lambda_j} \leq 0.$$
But all terms here are non-negative (since $|\lambda_j| \leq 1$), so this forces a contradiction unless $c_2=0$ (and hence also $c_1=0$) and $\lambda_j = 1$ for all $j$.
\end{remark}

As mentioned in Example \ref{ex2}, the most delicate regime occurs when $a = 1 - O(1/n)$.  It is convenient to introduce the normalized version
\begin{equation}\label{alpha-def} 
\alpha := \frac{n-1}{2} (1-a^2),
\end{equation}
of $a$, thus $0 < \alpha < \frac{n-1}{2}$, and the case $a = 1 - O(1/n)$ corresponds to $\alpha = O(1)$.  Informally, $\alpha$ measures how close $a$ is to $1$ (at the scale of $O(1/n)$).

A key role in the argument will be played by the mean
\begin{equation}\label{xy-def}
 x + iy := \frac{1}{n-1} \sum_{j=1}^{n-1} q_j
\end{equation}
of the $q_j$, particularly the real part $x$.  As the $q_j$ all lie in the unit disk, the mean $x+iy$ does also, so that
$$ -1 \leq x \leq 1$$
and
\begin{equation}\label{ybound}
|y| \leq \sqrt{1-x^2}.
\end{equation}
On the other hand, in the example in Example \ref{ex}, $x$ is equal to the extremal value of $1$, and $y=0$.  In Example \ref{ex2}, we have $x = 1 - O(1/n)$ (and $y = O(1/n)$).

It will be convenient to work with the quadratic polynomial
\begin{equation}\label{beta-poly}
    \beta(t) := 1 - 2atx + a^2 t^2 = 1 - x^2 + (x - at)^2
\end{equation}
with a particular emphasis on the value at $t=1$:
\begin{equation}\label{beta-def} 
    \beta(1) = 1 - 2ax + a^2 = (1-a)^2 + 2a(1-x) = 1 - x^2 + (x-a)^2.
\end{equation}
One should primarily think of $\beta(1)$ as a measure of how close $x$ is to $1$. Clearly we have
$$ \beta(t) > 0$$
for all $0 \leq t \leq 1$ (note that $at$ is strictly less than $1$).

The arguments will revolve around the relationship between $\alpha$ and $\beta(1)$.  Specifically, we will establish the following two inequalities.  The first inequality, which will be derived easily from the polar identity, is exact:

\begin{proposition}[Polar inequality]\label{polar}  We have 
\begin{equation}\label{lt} 
1 < \int_0^1 \exp ( \alpha (-1 + (2-\beta(1)) t))\ dt.
\end{equation}
\end{proposition}

The proof will be given below the fold.  One can weaken this inequality slightly to obtain a simpler one:

\begin{corollary}[Simplified polar inequality]\label{polar-corollary}  We have
\begin{equation}\label{beta-bound}
    0 < \beta(1) < \frac{\alpha}{3+\alpha} < 1.
\end{equation}
In particular, since $1 - \beta(1) = 2a(x-a/2)$, one has
\begin{equation}\label{xa} 
    x > a/2 > 0.
\end{equation}
\end{corollary}

The implication of \eqref{beta-bound} from \eqref{lt} can be illustrated by a numerical plot:

???

As the picture suggests, the implication is rather delicate; a formal derivation, provided by an AI, is given below the fold.


The second inequality, which is more difficult, is mostly deduced from the origin and centroid identities (with a minor assist from the polar inequality also playing a minor role to manage error terms):

\begin{proposition}[Origin inequality]\label{origin} 
If $n \geq 5$, then
\begin{equation}\label{origin-exact}
 2 \alpha + ax \leq \frac{1-x^2}{2(n-1)} + a^2 n(n-1) \int_0^1 t \beta(t)^{\frac{n-2}{2}}\ dt.
\end{equation}
\end{proposition}

With some computation, these inequalities give some further consequences:

\begin{corollary}\label{origin-cor}  Let $n \geq 5$.
\begin{itemize}
    \item [(i)]  We have $\alpha \leq 17$.
    \item [(ii)]  We have
\begin{equation}\label{1le}
     1 \leq \frac{\beta(1)}{2\alpha(1-\beta(1))} + \frac{\beta(1)}{4\alpha} + \frac{1}{2(n-1)} + \frac{\beta(1)}{4\alpha(n-1)}  
+ \frac{a^4 n (n-1) (n-2) \beta(1)}{4\alpha} \int_0^1 t^3 \beta(t)^{\frac{n-4}{2}}\ dt.
\end{equation}
\end{itemize}
\end{corollary}

We establish these bounds below the fold.

For fixed $\alpha$, the right-hand side of \eqref{1le} is monotone increasing in $\beta(1)$ (or equivalently, monotone decreasing in $x$).  In view of \eqref{beta-bound}, we can thus replace $\beta(1)$ by $\frac{\alpha}{3+\alpha}$ in this inequality, so that $ax = 1 - \frac{\alpha}{n-1} - \frac{\beta(1)}{2}$ is replaced by 
$$c(\alpha) := 1 - \frac{\alpha}{n-1} - \frac{\alpha}{2(3+\alpha)},$$ 
and $\beta(t) = 1 - 2axt + a^2t^2$ replaced by $1 - 2c(\alpha) t + a^2 t^2$. The inequality \eqref{1le} then becomes
\begin{equation}\label{stat} 1 \leq \frac{1}{6} + \frac{1}{4(3+\alpha)} + \frac{1}{2(n-1)} + \frac{1}{4(n-1)(3+\alpha)} + \frac{a^4 n (n-1) (n-2)}{4(3+\alpha)} \int_0^1 t^3 (1 - 2c(\alpha) t + a^2 t^2)^{\frac{n-4}{2}}\ dt.
\end{equation}
We also note that the upper bound $c(\alpha) \leq xa$ forces the constraint 
\begin{equation}\label{ca}
c(\alpha)^2 \leq a^2 = 1 - \frac{2\alpha}{n-1}.
\end{equation}
This prevents $\alpha$ from getting too close to the upper limit $\frac{n-1}{2}$ (or $a$ getting too close to zero).

We can now eliminate all large degrees, e.g., $n > 200$, as follows.  The quadratic $1 - 2c(\alpha) t + a^2 t^2$ attains its minimum at $t =c(\alpha)/a^2$.  For $t \leq c(\alpha)/a^2$ we have
$$ 1 - 2c(\alpha) t + a^2 t^2 \leq 1 - c(\alpha) t \leq \exp( - c(\alpha) t)$$
while for $c(\alpha) / a^2 < t \leq 1$ (if this region is non-vacuous) we can bound the quadratic by its value $\frac{\alpha}{3+\alpha}$ at $t=1$.  Thus
$$ \int_0^1 t^3 (1 - 2c(\alpha) t + a^2 t^2)^{\frac{n-4}{2}}\ dt
\leq \int_0^\infty t^3 \exp( - \frac{(n-4)}{2} c(\alpha) t)\ dt + \int_0^1 t^3 (\frac{\alpha}{3+\alpha})^{\frac{n-4}{2}}\ dt.$$
Evaluating these expressions, we arrive at
$$ 1 \leq \frac{1}{6} + \frac{1}{4(3+\alpha)} + \frac{1}{2(n-1)} + \frac{1}{4(n-1)(3+\alpha)} + \frac{24 a^4 n (n-1) (n-2)}{(3+\alpha) (n-4)^4 c(\alpha)^4}
+ \frac{a^4 n (n-1) (n-2)}{16(3+\alpha)} (\frac{\alpha}{3+\alpha})^{\frac{n-4}{2}}.$$
Since $\alpha \leq 17$, we have $\frac{\alpha}{3+\alpha} \leq \frac{17}{20}$.  
Next, we claim that $(3+\alpha) c(\alpha)^4 \geq 1$.  As $c(\alpha)$ is monotone increasing in $n$, it suffices to do this when $n=201$.  Here one can directly compute that
$$ \frac{d}{d\alpha} ((3+\alpha) c(\alpha)^4) = - c(\alpha)^3 \frac{5(3+\alpha)^2 - 103(3+\alpha) + 900}{200(3+\alpha)} < 0$$
since the discriminant $-7391$ of the the numerator is negative discriminant, we conclude that
$$ (3+\alpha) c(\alpha)^4 \geq (3+17) c(17)^4 \geq 1.1529\dots > 1$$
as desired.

Dropping some $\alpha$ and $a$ terms, we conclude that
$$ 1 \leq \frac{1}{6} + \frac{1}{12} + \frac{7}{12(n-1)} + \frac{24 n (n-1) (n-2)}{(n-4)^4}
+ \frac{n (n-1) (n-2)}{48} (\frac{17}{20})^{\frac{n-4}{2}}.$$
Every term on the right-hand side can be seen to be decreasing in $n$ for $n \geq 201$.  Thus the right-hand side can be bounded by
$$ \frac{1}{6} + \frac{1}{12} + \frac{7}{12 \times 200} + \frac{24 \cdot 201 \cdot 200 \cdot 199}{197^4}
+ \frac{201 \cdot 200 \cdot 199}{48} (\frac{17}{20})^{\frac{201-4}{2}} \leq 0.399,$$
giving the desired contradiction.

The remaining range to handle is when
$$ 5 \leq n \leq 200; 0 \leq \alpha \leq 17; \quad c(\alpha)^2 \leq 1 - \frac{2\alpha}{n-1}.$$
It turns out that \eqref{stat} remains infeasible in this range.  This can be illustrated numerically without much difficulty: see \href{???}{this applet}.  A rigorous certificate of this in interval arithmetic or Lean should be possible, but I have not yet attempted to do so. 


\more

\section{The polar inequality}

We begin with a proof of Proposition \ref{polar} and Corollary \ref{polar-corollary}.

As is well known, the M\"obius transform $z \mapsto \frac{a-z}{1-az}$ maps the closed unit disk to itself.  In particular, we have
$$ |\frac{1-az_j}{a-z_j}| \geq 1$$
for all of the zeroes $z_j$. Inserting this into the polar identity in Lemma \ref{identities}(ii) and using the triangle inequality, we conclude the lower bound
$$ \int_0^1 \prod_{j=1}^{n-1} |t (1-a^2) q_j + a|\ dt \geq 1.$$
We now convert this bound to a bound involving the quantity $x$ in \eqref{xy-def}., Fom the arithmetic mean-geometric mean inequality we have
\begin{equation}\label{amgm} \prod_{j=1}^{n-1} |t (1-a^2) q_j + a| \leq (\frac{1}{n-1} \sum_{j=1}^{n-1} |t (1-a^2) q_j + a|^2)^{\frac{n-1}{2}}
\end{equation}
and from \eqref{xy-def} we have
$$\sum_{j=1}^{n-1} |t (1-a^2) q_j + a|^2 = (n-1) a^2 + 2 (n-1) a t (1-a^2) x + t^2 (1-a^2)^2 \sum_{j=1}^{n-1} |q_j|^2.$$
Since $|q_j|^2 \leq 1$, we thus have
\begin{equation}\label{quad}
 \frac{1}{n-1} \sum_{j=1}^{n-1} |t (1-a^2) q_j + a|^2 \leq a^2 + 2 a t x (1-a^2) + t^2 (1-a^2)^2
\end{equation}
giving the inequality
\begin{equation}\label{1Q}
1 \leq \int_0^1 (a^2 + 2 a t x (1-a^2) + t^2 (1-a^2)^2)^{\frac{n-1}{2}}\ dt.
\end{equation}
We now use some slightly crude estimates to simplify the right-hand side of \eqref{1Q}.
Bounding $t^2$ by $t$ and using the quantities $\alpha,\beta(1)$ from \eqref{alpha-def}, \eqref{beta-def}, we observe that
\begin{equation}\label{at}
 a^2 + 2 a t x (1-a^2) + t^2 (1-a^2)^2 \leq
1 + \frac{2}{n-1} \alpha (-1 + (2-\beta(1)) t).
\end{equation}
Using the basic inequality $1 + y \leq \exp(y)$, we thus have
$$ (a^2 + 2 a t x (1-a^2) + t^2 (1-a^2)^2)^{\frac{n-1}{2}} \leq \exp ( \alpha (-1 + (2-\beta(1)) t)),$$
with strict inequality for $t < 1$. From \eqref{1Q} we conclude \eqref{lt}, which establishes Proposition \ref{polar}.

Now we prove Corollary \ref{polar-corollary}.  Here we use an AI-generated argument. One can directly calculate
$$ \int_0^1 \exp ( \alpha (-1 + (2-\beta(1)) t))\ dt = e^{-u} \frac{\sinh h}{h}$$
where $u = \alpha \beta(1)/2$ and $h = \alpha (2-\beta(1))/2$.  If we can show that
\begin{equation}\label{lsh}
 \log \frac{\sinh h}{h} \leq \sqrt{h^2+9} - 3
\end{equation}
for all $h > 0$, then taking logarithms in \eqref{lt} yields
$$ u < \log \frac{\sinh h}{h} \leq \sqrt{h^2+9} - 3$$
from which Corollary \ref{polar-corollary} follows by routine algebra.

Both sides of \eqref{lsh} vanish at $h=0$.  Taking derivatives, it suffices to show that
$$ \coth h - \frac{1}{h} \leq \frac{h}{\sqrt{h^2+9}}$$
which rearranges to
$$ h^4 \sinh^2 h - (h^2+9) (h \cosh h - \sinh h)^2 \geq 0.$$
But the left-hand side has a Taylor expansion
$$ \sum_{k=4}^\infty \frac{2^{2k-3} (2k-1) (2k-6)^2}{(2k)!} h^{2k},$$
giving the claim.



\section{The origin inequality}

Now we turn to the proof of Proposition \ref{origin}, which is more difficult and revolves around an analysis of the function $F$ defined in \eqref{F-def}.  We begin with a heuristic analysis.  Inserting the approximation $1 - \eps \approx \exp(-\eps)$ for small $\eps$ into \eqref{F-def} and using \eqref{xy-def}, we are led to the approximation
\begin{equation}\label{ft}
 F(t) \approx \exp( - (n-1) a t (x+iy) ),
\end{equation}
at least when $t$ is small (which turns out to be the dominant regime in applications).  This suggests a relation
\begin{equation}\label{1fq}
 \int_0^1 F(t)\ dt \approx \frac{1 - F(1)}{(n-1) a (x+iy)}
\end{equation}
between the two expressions involving $F$ in the origin identities in Lemma \ref{identities}.  Substituting in this approximation, we obtain some (slightly complicated) approximation for the sum $\sum_{j=1}^{n-1} \frac{1}{z_j}$ in terms of $a$, $n$, $x$, $y$, and the product $J := \prod_{j=1}^{n-1} z_j q_j$.

As $z_j, q_j$ lie in the closed unit disk, the product $J$ does also.  However, past experience with the Sendov conjecture has taught us that the worst cases tend to be when $z_j, q_j$ lie very close to the boundary of the disk, so that $|J|$ is close to one.  For instance, in Example \ref{ex} all the $z_j$ and $q_j$ lie on the unit circle, and $J = (-1)^{n+1}$.  See ??? for other places where this heuristic is noted.  To simplify the discussion, let us assume for now that $|J|$ is exactly one, so that $z_j, q_j$ all lie on the unit circle.  This leads in particular to the inversion identities
\begin{equation}\label{inverse}
     \frac{1}{z_j} = \overline{z_j}; \quad \frac{1}{q_j} = \overline{q_j}.
\end{equation}
The centroid identity in Lemma \ref{identities}(i) relate the sum of the $z_j$ with the sum of the $1/q_j$.  Using \eqref{inverse}, this gives a similar identity relating the sum of the $1/z_j$ with the sum of the $q_j$.  The latter sum is of course just $(n-1) (x+iy)$.  This combines well with the previous approximation, thus giving an approximate identity relating $x$, $y$ to $J$, $a$, and $n$.  As it turns out, the roles of $y$ and $J$ are minor and can be quickly eliminated for the purposes of obtaining useful bounds, leading eventually to the relation in Proposition \ref{origin}.

We turn to the details.  To make the approximation \eqref{1fq} more precise, we note that $F(0)=1$, and hence by the fundamental theorem of calculus
$$ 1 = F(1) - \int_0^1 F'(t)\ dt.$$
The heuristic \eqref{ft} predicts that $F'(t) \approx -(n-1) a(x+iy) F(t)$, which would give \eqref{1fq}.  If we actually differentiate \eqref{F-def} carefully, we obtain the exact identity
$$ F'(t) = - (n-1) a(x+iy) F(t) - \sum_{j=1}^{n-1} a^2 t q^2_j \prod_{k \neq j} (1 - atq_k).$$
Bounding $|q^2_j| \leq 1$, we write this 
$$ F'(t) = - (n-1) a(x+iy) F(t) + O_{\leq} ( a^2 t \sum_{j=1}^{n-1} \prod_{k \neq j} |1 - atq_k| ).$$

When faced with a similar expression in \eqref{amgm}, we used the arithmetic mean-geometric mean inequality.  Here, the analogous tool is \href{https://en.wikipedia.org/wiki/Maclaurin\%27s_inequality}{Maclaurin's inequality}, which gives
$$
\frac{1}{n-1} \sum_{j=1}^{n-1} \prod_{k \neq j} |1 - atq_k| \leq (\frac{1}{n-1} \sum_{j=1}^{n-1} |1 - at q_j|)^{n-2},$$
and hence by Cauchy--Schwarz
$$
\frac{1}{n-1} \sum_{j=1}^{n-1} \prod_{k \neq j} |1 - atq_k| \leq (\frac{1}{n-1} \sum_{j=1}^{n-1} |1 - at q_j|^2)^{\frac{n-2}{2}}.$$
Repeating the calculations used to show \eqref{quad}, we have
$$ \frac{1}{n-1} \sum_{j=1}^{n-1} |1 - at q_j|^2 \leq 1 - 2 a x t + a^2 t^2 = \beta(t)$$
and so we obtain the bound
$$ F'(t) = - (n-1) a(x+iy) F(t) + O_{\leq} ( a^2 t (n-1) \beta(t)^{\frac{n-2}{2}} ).$$
Integrating this, we obtain a rigorous analogue of \eqref{1fq},
$$ 1 - F(1) = (n-1) a(x+iy) \int_0^1 F(t)\ dt + O_{\leq} ( \int_0^1 a^2 t (n-1) \beta(t)^{\frac{n-2}{2}}\ dt )$$
and thus by the triangle inequality
\begin{equation}\label{tri}
 1 \leq |F(1) + (n-1) a(x+iy) \int_0^1 F(t)\ dt| + \int_0^1 a^2 t (n-1) \beta(t)^{\frac{n-2}{2}}\ dt.
\end{equation}
From the first and second origin identities we have
\begin{equation}\label{f1aq}
 F(1) + (n-1) a(x+iy) \int_0^1 F(t)\ dt = \frac{(-1)^{n-1} J}{n} 
 \left(1 + a \sum_{j=1}^{n-1} \frac{1}{z_j} + (n-1) a(x+iy) \right).
\end{equation}

The next step is thus to estimate $\sum_{j=1}^{n-1} \frac{1}{z_j}$.  When $|J|=1$, then all the $z_j, q_j$ were on the unit circle and we could use \eqref{inverse} (and the centroid identity) to proceed.  Now, we are no longer assuming $|J|$ to equal $1$, but we can still adapt the previous arguments with a loss proportional to $1 - |J|^2$.  The key lemma is

\begin{lemma}[Defect lemma]  Let $w_1,\dots,w_N$ be some points in the closed unit disk.  Then
$$ \prod_{j=1}^N |w_j| \times \sum_{j=1}^N |\frac{1}{w_j} - \overline{w_j}| \leq 1 - \prod_{j=1}^N |w_j|^2.$$
\end{lemma}

\begin{proof} By a limiting argument we may assume that none of the $w_j$ vanish.  If we write $|w_j| = e^{-a_j}$ for some $a_j \geq 0$, then we can calculate that
    $$ |\frac{1}{w_j} - \overline{w_j}| = 2 \sinh a_j$$
    and
$$ 1 - \prod_{j=1}^N |w_j|^2 = \prod_{j=1}^N |w_j| \times 2 \sinh \sum_{j=1}^N a_j.$$
Thus the desired inequality reduces to the superadditivity property
$$ \sinh \sum_{j=1}^N a_j \geq \sum_{j=1}^N \sinh a_j.$$
But from the sinh addition formula $\sinh(a+b) = \sinh a \cosh b + \sinh b \cosh a$ we have
$$ \sinh(a+b) \geq \sinh a + \sinh b$$
for all non-negative $a,b$ (this also follows from the convex nature of $\sinh$ together with $\sinh 0 = 0$), and the claim follows by induction.
\end{proof}

From taking complex conjugates of the centroid identity and performing some algebra, we have
$$
J \sum_{j=1}^{n-1} \left(\frac{n-1}{n} \overline{z_j} + \frac{1}{\overline{q_j}}\right) = a J \frac{(n-1)^2}{n}.$$
Using the defect lemma (applied to the points $z_1,\dots,z_{n-1},\overline{q_1},\dots,\overline{q_{n-1}}$) and the triangle inequality we conclude that
$$
J \sum_{j=1}^{n-1} \left(\frac{n-1}{n} \frac{1}{z_j} + q_j\right) = a J \frac{(n-1)^2}{n}
+ O_{\leq}( 1 - |J|^2 ),$$
where as before we are removing singularities when some of the $z_j$ vanish.
and hence after applying \eqref{xy-def} and some algebraic manipulation, we arrive at
$$
J \sum_{j=1}^{n-1} \frac{1}{z_j} = a J (n-1) - n J (x+iy)
+ O_{\leq}( \frac{n}{n-1} (1 - |J|^2) )$$
Substituting this back into \eqref{f1aq}, we conclude that
$$\left|F(1) + (n-1) a(x+iy) \int_0^1 F(t)\ dt\right| = 
\left| \frac{|J|}{n} + \frac{a^2 |J|(n-1)}{n} - a|J| (x+iy) + \frac{a (n-1) (x+iy) |J|}{n} \right|
+ O_{\leq}( \frac{a}{n-1} (1-|J|^2) )$$
and hence after some algebra and the triangle inequality
$$\left|F(1) + (n-1) a(x+iy) \int_0^1 F(t)\ dt\right| \leq 
|J| \left| \frac{a^2 (n-1)}{n} + \frac{1}{n} - \frac{a(x+iy)}{n} \right|
+ \frac{a}{n-1} (1-|J|^2).$$
Inserting this into \eqref{tri}, we obtain
\begin{equation}\label{tri-2}
 1 \leq |J| \left| \frac{a^2 (n-1)}{n} + \frac{1}{n} - \frac{a(x+iy)}{n} \right| + \frac{a}{n-1} (1-|J|^2) + \int_0^1 a^2 t (n-1) \beta(t)^{\frac{n-2}{2}}\ dt.
\end{equation}
We can simplify \eqref{tri-2} by reducing to the $|J|=1$ case.  Indeed, we shall show that
\begin{equation}\label{grow}
    \left| \frac{a^2 (n-1)}{n} + \frac{1}{n} - \frac{a(x+iy)}{n} \right|  \geq 2 \frac{a}{n-1}
\end{equation}
which implies that the right-hand side of \eqref{tri-2} is non-decreasing in $|J|$ in the range $0 \leq |J| \leq 1$.  Thus we may replace $|J|$ by $1$ in \eqref{tri-2} to conclude that
\begin{equation}\label{tri-3}
 1 \leq \left| \frac{a^2 (n-1)}{n} + \frac{1}{n} - \frac{a(x+iy)}{n} \right| + \int_0^1 a^2 t (n-1) \beta(t)^{\frac{n-2}{2}}\ dt.
\end{equation}
Let us now verify \eqref{grow}.  Using $|x+iy| \leq 1$ and the triangle inequality, we can lower bound
$$     \left| \frac{a^2 (n-1)}{n} + \frac{1}{n} - \frac{a(x+iy)}{n} \right|  
\geq  \frac{a^2 (n-1)}{n} + \frac{1}{n} - \frac{a}{n}.$$
Inserting this into \eqref{grow} and clearing denominators, we reduce after some algebra to
$$ (n-1)^2 a^2 - (3n-1) a + n-1 > 0.$$
But as a quadratic polynomial in $a$, the left-hand side has discriminant
$(3n-1)^2 - 4(n-1)^3$, which one can check to be negative for sufficiently large $n$ (in fact $n \geq 5$ suffices), giving the claim \eqref{grow}.

Next we eliminate the role of the imaginary term $iy$.
 Observe for any complex number $s+it$ with positive real part that
$$ |s+it| \leq s + \frac{t^2}{2s}$$
as can be seen by squaring both sides.  The expression $\frac{a^2 (n-1)}{n} + \frac{1}{n} - \frac{a(x+iy)}{n}$ has real part
$$ \frac{a^2 (n-1)}{n} + \frac{1-ax}{n}$$
which lies between $\frac{a^2 (n-1)}{n}$ and $1$ (in particular, it is positive), and imaginary part of magnitude at most
$$ \frac{|a|}{n} (1 - x^2)^{1/2}$$
by \eqref{ybound}. We conclude that
$$ \left| \frac{a^2 (n-1)}{n} + \frac{1}{n} - \frac{a(x+iy)}{n} \right|  
\leq \frac{a^2 (n-1)}{n} + \frac{1-ax}{n} + \frac{1-x^2}{2n(n-1)}.$$
The right-hand side can be rearranged using the quantity $\alpha$ from \eqref{alpha-def} as
$$ 1 - \frac{2\alpha}{n}  - \frac{ax}{n} + \frac{1-x^2}{2n(n-1)},$$
so the bound \eqref{tri-3} gives \eqref{origin-exact}.

\subsection{Upper bound on $\alpha$}

Now we can prove Corollary \ref{origin-cor}(i). Suppose for contradiction that $\alpha > 17$; since $\alpha \leq \frac{n-1}{2}$, this implies that $n \geq 36$.
Crudely discarding the $ax$ term in \eqref{origin-exact} and bounding $\frac{1-x^2}{2(n-1)}$ by $\frac{1}{2(n-1)}$, we have
$$ 2 \alpha \leq \frac{1}{2(n-1)} + a^2 n(n-1) \int_0^1 t \beta(t)^{\frac{n-2}{2}}\ dt.$$
The quadratic polynomial $\beta(t)$ equals $1$ at $t=0$ and attains its minimum at $t = x/a$ with value $1-x^2$.  By convexity, we thus have
\begin{equation}\label{ax}
 \beta(t) \leq 1 - axt \leq \exp(-axt) 
\end{equation}
for $0 \leq t \leq x/a$ and
\begin{equation}\label{tax}
 \beta(t) \leq \beta(1)
\end{equation}
for $x/a \leq t \leq 1$ (this latter statement is vacuous if $x/a > 1$).  Since $\int_0^1 t\ dt = \frac{1}{2}$,  We can therefore crudely bound
$$
\int_0^1 t (1 - 2 a x t + a^2 t^2)^{\frac{n-2}{2}}\ dt \leq \int_0^\infty t \exp( - \frac{n-2}{2} axt)\ dt + \frac{1}{2} \beta(1)^{\frac{n-2}{2}} $$
$$ = \frac{4}{(n-2)^2 a^2 x^2} + \frac{1}{2} \beta(1)^{\frac{n-2}{2}},$$
and hence
$$ 2 \alpha \leq \frac{1}{2(n-1)} + \frac{4n(n-1)}{(n-2)^2 x^2} + \frac{a^2 n(n-1)}{2} \beta(1)^{\frac{n-2}{2}}.$$
From evaluating the integral in \eqref{lt} one has
$$1 < \frac{e^{\alpha (1 - \beta(1))}- e^{-\alpha}}{\alpha (2 - \beta(1))}$$
thus
$$ e^{\alpha (1 - \beta(1))} > \alpha (2 - \beta(1)) + e^{-\alpha} \geq \alpha,$$
so on taking logarithms we have
\begin{equation}\label{beta-lower}
1 - \beta(1) > \frac{\log \alpha}{\alpha}.
\end{equation}
From \eqref{beta-def} we have $x^2 \geq 1-\beta(1)$, thus by \eqref{beta-lower} one has
$$ \frac{1}{x^2} \leq \frac{\alpha}{\log \alpha}.$$
From another application of \eqref{beta-lower} one has
$$ \beta(1) \leq \exp( - (1 - \beta(1)) ) \leq \alpha^{-1/\alpha}.$$
We conclude that
$$ 2 \alpha \leq \frac{1}{2(n-1)} + \frac{4n(n-1)}{(n-2)^2 \log \alpha} \alpha + \frac{a^2 n(n-1)}{2} \alpha^{-\frac{n-2}{2\alpha}}.$$
It is now convenient to introduce the quantity $u := \frac{a^2}{1-a^2}$, thus $0 < u < \infty$ with
$$ a^2 = \frac{u}{1+u}$$
and
$$ \frac{n-2}{2 \alpha} = 1 + u - \frac{1}{2\alpha}.$$
Inserting these bounds and dividing by $\alpha$, we conclude
$$ 2 \leq \frac{1}{2(n-1) \alpha} + \frac{4n(n-1)}{(n-2)^2 \log \alpha} + \frac{u}{2(1+u) \alpha^2} n(n-1) \alpha^{-u} \alpha^{\frac{1}{2\alpha}}.$$
Since $\alpha = \frac{n-1}{2(1+u)}$, we obtain
$$ 2 \leq \frac{1}{2(n-1) \alpha} + \frac{4n(n-1)}{(n-2)^2 \log \alpha} + \frac{2n}{n-1} u (1+u) \alpha^{-u} \alpha^{\frac{1}{2\alpha}}.$$
Since $n \geq 36$ and $\alpha \geq 17$, we conclude that
$$ 2 \leq \frac{1}{2 \times 35 \times 17} + \frac{4 \times 36 \times 35}{(34)^2 \log 17} + \frac{2 \times 36}{35} u (1+u) 17^{-u} 17^{\frac{1}{2 \times 17}}.$$
Routine calculus shows that $u(1+u) 17^{-u}$ has a maximum of at most $0.1825$, and that the right-hand side here is at most $1.948$, giving the required contradiction.  This proves Corollary \ref{origin-cor}(i).


\subsection{A simplified estimate}

Now we show Corollary \ref{origin-cor}(ii).
Note from \eqref{alpha-def} that
\begin{equation}\label{1a}
 1-a^2 = \frac{2\alpha}{n-1}
\end{equation}
while from \eqref{beta-def} we have
\begin{equation}\label{1x}
 1 - x^2 \leq \beta(1) 
\end{equation}
and hence also
\begin{equation}\label{1x'}
\frac{1}{x^2} \leq 1 + \frac{\beta(1)}{1-\beta(1)}.
\end{equation}
From \eqref{beta-def} we have
$$ \beta(t) = (1 - axt)^2 + a^2 t^2 (1-x^2).$$
By the mean value theorem (noting that $1-axt$ is non-negative) we thus have
$$ \beta(t)^{\frac{n-2}{2}} \leq ((1 - axt)^2)^{\frac{n-2}{2}}
+ a^2 t^2 (1-x^2) \frac{n-2}{2} \beta(t)^{\frac{n-4}{2}}.$$
From the standard beta function identity
$$ \int_0^{1/ax} t (1 - axt)^{n-2}\ dt = \frac{1}{a^2 x^2 n(n-1)}$$
(and the fact that $1/ax \geq 1$) we can thus replace \eqref{origin-exact} by
$$ 2\alpha + ax \leq \frac{1-x^2}{2(n-1)} + \frac{1}{x^2}
+ \frac{a^4 n (n-1) (n-2) (1-x^2)}{2} \int_0^1 t^3 \beta(t)^{\frac{n-4}{2}}\ dt.$$
From \eqref{beta-def} we have
$$ax = 1 - \frac{\beta(1)}{2} - \frac{1-a^2}{2}$$
Thus by \eqref{1a}, \eqref{1x}, \eqref{1x'} 
$$ 2\alpha \leq \frac{\beta(1)}{1-\beta(1)} + \frac{\beta(1)}{2} + \frac{\alpha}{n-1} + \frac{\beta(1)}{2(n-1)}  
+ \frac{a^4 n (n-1) (n-2) \beta(1)}{2} \int_0^1 t^3 \beta(t)^{\frac{n-4}{2}}\ dt.$$
Dividing by the positive quantity $2\alpha$ gives the claim.


\end{document}
